% titel: Blatt 2, A1: Position nach einem Zufallswalk

% Teil I: Werkzeuge, kein Lösungsweg.

% Grundidee: Position = R-L = 2R-n. Parität von Position und n stimmen
% überein. Symmetrie R<->L plus eventueller Gleichstand bei 0.

Ein Zufallswalk ist das einfachste Modell eines ``irrenden`` Prozesses:
pro Schritt $\pm1$ mit Wahrscheinlichkeit je $\tfrac12$. Zwei Fragen sind
dabei zentral: \emph{wo} kann man nach $n$ Schritten überhaupt stehen
(Parität) und \emph{wie wahrscheinlich} ist es, rechts zu stehen
(Symmetrie).

\begin{werkzeugkasten}[Zentrale Beziehung: Position = $2R-n$]
Bezeichne $R$ die Zahl der Rechtsschritte, $L$ die der Linksschritte und
$R+L=n$. Die Position nach $n$ Schritten ist
\begin{equation}
  \text{Position} = R-L = R-(n-R)=2R-n .
\end{equation}
\textbf{Parität:} $R$ ganzzahlig, also ist $2R$ gerade; die Position
$2R-n$ hat damit \emph{dieselbe Parität wie $n$}. Bei geradem $n$ sind nur
gerade Positionen, bei ungeradem $n$ nur ungerade Positionen erreichbar:
\begin{equation}
  \text{erreichbar}=\{-n,-n+2,\dots,n\}.
\end{equation}
\end{werkzeugkasten}

\begin{werkzeugkasten}[Wahrscheinlichkeit, rechts zu stehen]
$R$ ist binomialverteilt, $R\sim\mathrm{Bin}(n,\tfrac12)$, mit
$\mathrm P(R=r)=\binom{n}{r}2^{-n}$.
\begin{itemize}
  \item $n$ \emph{ungerade}: Die Verteilung ist symmetrisch um $0$
        ($R\leftrightarrow L$), es gibt keinen Gleichstand, also
        $\mathrm P(\text{rechts})=\mathrm P(\text{links})=\tfrac12$.
  \item $n$ \emph{gerade}: Position $0$ ist erreichbar, mit
        $\mathrm P(0)=\binom{n}{n/2}2^{-n}$; aus der Symmetrie und der
        Normierung ($P($rechts$)+P($links$)+P(0)=1$) folgt
        $\mathrm P(\text{rechts})=\frac12\bigl(1-\mathrm P(0)\bigr)$.
\end{itemize}
``Kann das Individuum zur Startposition zurückkehren?`` lautet genau für
\emph{gerades} $n$ mit ``ja``.
\end{werkzeugkasten}

\begin{achtung}[Darauf achten]
\begin{itemize}
  \item Parität: Position und $n$ haben dieselbe Parität --- bei geradem
        $n$ nur gerade Positionen, $0$ eingeschlossen.
  \item Bei geradem $n$ ist $P(\text{rechts})<\tfrac12$, weil Masse bei
        $0$ liegt; bei ungeradem $n$ exakt $\tfrac12$.
  \item Nicht jede Position \emph{gleich} wahrscheinlich: mittlere
        Positionen (nahe $R\approx n/2$) sind wahrscheinlicher
        (Binomialhügel).
\end{itemize}
\end{achtung}
